% =============================================================================
% appendix.tex — Приложения
% =============================================================================

\appendix
\renewcommand{\thesection}{\Asbuk{section}}

\section*{ПРИЛОЖЕНИЯ}
\addcontentsline{toc}{section}{ПРИЛОЖЕНИЯ}

\subsection*{Приложение А. Схема базы данных}
\addcontentsline{toc}{subsection}{Приложение А. Схема базы данных}

\todo{ER-диаграмма (PostgreSQL 16 + pgvector, 23 SQL-миграции Liquibase):

Основные таблицы:
\begin{enumerate}
  \item \textbf{application} — реестр приложений: id, key (unique), name, description, repo\_url, repo\_provider (github|gitlab|bitbucket), repo\_owner, repo\_name, monorepo\_path, default\_branch, last\_commit\_sha, last\_indexed\_at, last\_index\_status (success|failed|partial|running), owners (JSONB: List<OwnerRef>), contacts (JSONB: List<ContactRef>), tags (text[]), languages (text[]), embedding\_model, embedding\_dim (default 1024), ann\_index\_params (JSONB: method=ivfflat, lists=100), rag\_prefs (JSONB: priority, reranker), retention\_days, pii\_policy (JSONB), metadata (JSONB)
  \item \textbf{node} — узлы графа: id, application\_id (FK), fqn, kind (NodeKind enum), lang (Lang enum), parent\_id (FK, самореференция), file\_path, line\_start, line\_end, source\_code (TEXT), doc\_comment (TEXT), signature, code\_hash, meta (JSONB)
  \item \textbf{edge} — рёбра графа: src\_id (FK), dst\_id (FK), kind (EdgeKind enum) — составной PK; evidence (TEXT), confidence (FLOAT), relation\_strength (FLOAT)
  \item \textbf{chunk} — RAG-чанки: id (UUID), application\_id (FK), node\_id (FK), source, kind, content (TEXT), content\_hash, embedding (vector(1024)), embed\_model, metadata (JSONB)
  \item \textbf{node\_doc} — каноническая документация: node\_id (PK), locale, summary, details, examples, quality (JSONB), source\_kind, model\_meta (JSONB)
  \item \textbf{library} — внешние JAR: id, coordinate (unique), group\_id, artifact\_id, version, kind (framework|library|language|company|external)
  \item \textbf{library\_node} — узлы из JAR: id, library\_id (FK), fqn, kind, parent\_id (FK)
  \item \textbf{node\_library\_node\_edge} — связи между узлами кода и узлами из библиотек
  \item \textbf{ingest\_run} — запуски индексации: id, application\_id (FK), status, started\_at, finished\_at, commit\_sha
  \item \textbf{ingest\_step} — шаги индексации: id, run\_id (FK), type (IngestStepType enum), status, started\_at, finished\_at
  \item \textbf{ingest\_event} — события: id, step\_id (FK), level (INFO|WARN|ERROR), message, created\_at
  \item \textbf{api\_key} — API-ключи: id, name, key\_hash (SHA-256), scopes (text[]), is\_active, expires\_at, last\_used\_at
  \item \textbf{audit\_log} — журнал аудита: id, user\_id, action, resource, http\_method, request\_body, ip\_address, created\_at
  \item \textbf{ai\_documents} — Spring AI pgvector store: id, embedding (vector), metadata (JSONB), content
  \item \textbf{ai\_chat\_memory} — Spring AI chat memory: id, conversation\_id, messages (JSONB)
  \item \textbf{synonym\_dictionary} — словарь синонимов: id, node\_id (FK), synonyms (text[]), status
\end{enumerate}

Индексы: trigram GIN на fqn (fuzzy search), tsvector на doc\_comment/summary/details (полнотекстовый поиск RU+EN), HNSW/IVFFlat на vector(1024) (ANN-поиск), BRIN на created\_at (временные запросы).

Расширения: pgvector, pg\_trgm, unaccent.}

\subsection*{Приложение Б. Диаграмма архитектуры системы}
\addcontentsline{toc}{subsection}{Приложение Б. Диаграмма архитектуры системы}

\todo{Диаграмма компонентов (11 модулей):

\begin{verbatim}
                     +-------------------+
                     |   Spring Boot App |
                     |   (src/main)      |
                     |  Controllers, API |
                     |  Security, Config |
                     +--------+----------+
                              |
          +-------------------+-------------------+
          |                                       |
   +------+------+                    +-----------+-----------+
   |   kernel/   |                    |      contexts/        |
   | domain(JPA) |                    |                       |
   | db (Repos)  |<---api/impl-------+ git    | graph        |
   | shared      |                    | chunking | embedding  |
   +---------+---+                    | ai     | rag          |
             |                        | library | postprocess |
             v                        +-----------+-----------+
   +-------------------+                        |
   | PostgreSQL 16     |                        v
   | + pgvector        |              +-------------------+
   +-------------------+              |   Ollama (LLM)    |
                                      |   qwen2.5-coder   |
   +-------------------+              |   qwen2.5-instruct|
   | Redis             |              |   bge-m3 (embed)  |
   | (rate limiting)   |              +-------------------+
   +-------------------+
\end{verbatim}

Docker Compose: 4 сервиса (db, ollama, evaluator, app).

Конвейер обработки (7 этапов):
GIT clone/pull → PSI-парсинг → GraphLinker (7 фаз) → LibraryBuilder (ASM) → ChunkBuildOrchestrator → EmbeddingStore → NodeDocGenerator (Coder→Talker→Digest)}

\subsection*{Приложение В. Перечисление типов узлов и связей}
\addcontentsline{toc}{subsection}{Приложение В. Перечисление типов узлов и связей}

\todo{Полная таблица NodeKind (24 типа):

\begin{center}
\begin{tabular}{|l|l|l|}
\hline
Тип & Категория & Описание \\
\hline
REPO & Код/структура & Корневой узел репозитория \\
MODULE & Код/структура & Gradle-модуль проекта \\
PACKAGE & Код/структура & Kotlin/Java пакет \\
CLASS & Код/структура & Класс (включая object) \\
INTERFACE & Код/структура & Интерфейс \\
ENUM & Код/структура & Перечисление \\
RECORD & Код/структура & data class \\
METHOD & Код/структура & Метод/функция \\
FIELD & Код/структура & Свойство/поле \\
EXCEPTION & Код/структура & Класс исключения \\
TEST & Код/структура & Тестовый класс/метод \\
MAPPER & Код/структура & Маппер (MapStruct и др.) \\
\hline
SERVICE & Сервисный слой & Логическая сервисная единица \\
ENDPOINT & Сервисный слой & REST/gRPC/WS входная точка \\
CLIENT & Сервисный слой & HTTP/WebClient/Feign/gRPC-клиент \\
TOPIC & Сервисный слой & Kafka/Rabbit тема/очередь \\
JOB & Сервисный слой & Scheduler/Worker/Batch \\
\hline
DB\_TABLE & Данные/инфра & Таблица БД \\
DB\_VIEW & Данные/инфра & Представление БД \\
DB\_QUERY & Данные/инфра & Запрос (MyBatis/JPA/SQL) \\
SCHEMA & Данные/инфра & Контракт данных (Avro/JSON Schema) \\
CONFIG & Данные/инфра & Параметр конфигурации \\
MIGRATION & Данные/инфра & Liquibase/Flyway миграция \\
ANNOTATION & Данные/инфра & Статическая аннотация \\
\hline
\end{tabular}
\end{center}

Полная таблица EdgeKind (24 типа):

\begin{center}
\begin{tabular}{|l|l|l|}
\hline
Тип & Категория & Пример \\
\hline
CONTAINS & Статика/ООП & REPO→MODULE, CLASS→METHOD \\
DEPENDS\_ON & Статика/ООП & CLASS→CLASS (импорт) \\
IMPLEMENTS & Статика/ООП & CLASS→INTERFACE \\
INHERITS & Статика/ООП & CLASS→CLASS (наследование) \\
EXTENDS & Статика/ООП & INTERFACE→INTERFACE \\
OVERRIDES & Статика/ООП & METHOD→METHOD \\
ANNOTATED\_WITH & Статика/ООП & CLASS→ANNOTATION \\
\hline
CALLS & Вызовы & (deprecated, заменён на CALLS\_CODE) \\
CALLS\_CODE & Вызовы & METHOD→METHOD \\
THROWS & Вызовы & METHOD→EXCEPTION \\
LOCKS & Вызовы & METHOD→RESOURCE \\
\hline
CALLS\_HTTP & Интеграции & CLIENT→ENDPOINT (HTTP) \\
CALLS\_CAMEL & Интеграции & SERVICE→ENDPOINT (Camel) \\
CALLS\_GRPC & Интеграции & CLIENT→ENDPOINT (gRPC) \\
PRODUCES & Интеграции & METHOD→TOPIC \\
CONSUMES & Интеграции & JOB→TOPIC \\
QUERIES & Интеграции & METHOD→DB\_QUERY \\
READS & Интеграции & SERVICE→DB\_TABLE \\
WRITES & Интеграции & SERVICE→DB\_TABLE \\
CONTRACTS\_WITH & Интеграции & ENDPOINT↔SCHEMA \\
CONFIGURES & Интеграции & CONFIG→SERVICE \\
CIRCUIT\_BREAKER\_TO & Отказоуст. & CLIENT→SERVICE \\
RETRIES\_TO & Отказоуст. & CLIENT→SERVICE \\
TIMEOUTS\_TO & Отказоуст. & CONFIG→SERVICE \\
\hline
\end{tabular}
\end{center}}

\subsection*{Приложение Г. Примеры сгенерированной документации}
\addcontentsline{toc}{subsection}{Приложение Г. Примеры сгенерированной документации}

\todo{Пример 1: Coder → Talker → Digest

Входной код (Kotlin):
\begin{verbatim}
@Service
class OrderService(
    private val orderRepository: OrderRepository,
    private val kafkaTemplate: KafkaTemplate<String, OrderEvent>,
) {
    fun createOrder(request: CreateOrderRequest): Order {
        val order = Order(/* ... */)
        orderRepository.save(order)
        kafkaTemplate.send("order-events", OrderCreatedEvent(order.id))
        return order
    }
}
\end{verbatim}

Результат Coder (docTech):
\begin{verbatim}
Purpose: Сервис управления заказами
Inputs: CreateOrderRequest (данные нового заказа)
Outputs: Order (созданный заказ)
SideEffects:
  - Сохранение в БД через OrderRepository
  - Отправка OrderCreatedEvent в Kafka topic "order-events"
Calls: OrderRepository.save, KafkaTemplate.send
\end{verbatim}

Результат Talker (docPublic):
<<Данный компонент отвечает за создание новых заказов в системе. При получении запроса на создание заказа, он сохраняет его в базу данных и отправляет уведомление о новом заказе в систему обмена сообщениями для дальнейшей обработки другими сервисами.>>

Пример 2: RAG-запрос

Запрос: <<Какие сервисы используют Kafka для обмена данными о заказах?>>

Шаги обработки:
1. EXTRACTION → entities: [OrderService, KafkaTemplate, order-events]
2. EXACT\_SEARCH → найден: OrderService (FQN)
3. GRAPH\_EXPANSION → связи: OrderService –PRODUCES→ order-events, PaymentService –CONSUMES→ order-events
4. VECTOR\_SEARCH → контекст: OrderCreatedEvent, PaymentProcessor
5. RERANKING → top-3 релевантных чанка

Ответ: <<В системе обмен данными о заказах через Kafka осуществляют два сервиса: OrderService (отправляет события OrderCreatedEvent в топик order-events) и PaymentService (получает эти события для обработки платежей).>>

Пример 3: скриншот визуализации графа (graph.html)
\begin{itemize}
  \item Single-App: граф с узлами CLASS, METHOD, ENDPOINT, TOPIC и рёбрами CALLS\_CODE, CALLS\_HTTP, PRODUCES
  \item Cross-App: два APPLICATION-узла, соединённых через ENDPOINT и TOPIC
\end{itemize}}

\subsection*{Приложение Д. Спецификация REST API}
\addcontentsline{toc}{subsection}{Приложение Д. Спецификация REST API}

\todo{Таблица эндпоинтов:

\textbf{Группа Ingestion} (/api/v1/ingest/):
\begin{center}
\begin{tabular}{|l|l|l|}
\hline
Метод & URL & Описание \\
\hline
POST & /start & Запуск индексации приложения \\
POST & /reindex & Полная переиндексация \\
GET & /status/\{appId\} & Статус текущей индексации \\
GET & /runs/\{appId\} & История запусков \\
GET & /events/\{runId\} & События запуска \\
\hline
\end{tabular}
\end{center}

\textbf{Группа Graph} (/api/v1/graph/):
\begin{center}
\begin{tabular}{|l|l|l|}
\hline
GET & /\{appId\} & Граф приложения (kinds, edgeKinds, limit, depth) \\
GET & /cross-app & Кросс-приложенческий граф \\
GET & /cross-app/applications & Список приложений \\
\hline
\end{tabular}
\end{center}

\textbf{Группа Search/RAG}:
\begin{center}
\begin{tabular}{|l|l|l|}
\hline
POST & /api/v1/rag/ask & RAG-запрос (query, sessionId, applicationId) \\
POST & /api/v1/embedding/search & Векторный поиск (query, topK) \\
POST & /api/v1/embedding/documents & Добавление документов \\
DELETE & /api/v1/embedding/clear-postprocess & Очистка \\
\hline
\end{tabular}
\end{center}

\textbf{Группа API Keys} (/api/v1/api-keys/):
\begin{center}
\begin{tabular}{|l|l|l|}
\hline
POST & / & Создание ключа \\
GET & / & Список ключей \\
DELETE & /\{id\} & Отзыв ключа \\
\hline
\end{tabular}
\end{center}

\textbf{Группа Monitoring}:
\begin{center}
\begin{tabular}{|l|l|l|}
\hline
GET & /actuator/health & Состояние (DB, Redis, Ollama) \\
GET & /actuator/prometheus & Метрики Prometheus \\
GET & /actuator/info & Информация о приложении \\
\hline
\end{tabular}
\end{center}

Rate limiting: Default=100 req/min, High-Load=30 req/min (/rag/ask, /embedding/search), Mutation=10 req/min (/ingest/start, /ingest/reindex).}
